\chapter{Conclusion and Future Work}
\label{chap:conclusion}

\section{Summary of Achievements}
\label{sec:summary}

This project set out to build a video sharing platform whose transcoding subsystem runs on a
serverless container and event-driven architecture, and to evaluate whether that choice
delivers the cost and scalability benefits it promises.

The functional platform was completed. Users can register, authenticate, upload videos directly
to Amazon S3 through pre-signed URLs, watch them with adaptive bitrate playback across three
renditions, control whether each video is public or private, manage their channel, search and
filter the catalogue, and interact through likes and comments.

The transcoding pipeline operates without manual intervention. An upload emits an S3 event that
is buffered in Amazon SQS and dispatched by an AWS Lambda function to AWS Batch, where a
container packaging FFmpeg produces the renditions, the master playlist, and a thumbnail before
terminating. Conditional database writes prevent a duplicate delivery from overwriting a completed
job, failed jobs are retried by AWS Batch and, when they exhaust their attempts, raise an alert
carrying the reason for the failure rather than terminating silently. Both mechanisms replaced
earlier designs that held only for an execution mode the deployed system does not use; the
reasoning is set out in Section~\ref{sec:pipeline-design} and Section~\ref{sec:challenges}.

Content protection was implemented in two layers. Origin Access Control keeps the storage
bucket entirely private, and CloudFront Signed Cookies ensure that private videos are rejected
at the edge even when the exact manifest URL is known. This closes a gap that would otherwise
have made the private visibility setting ineffective in practice, since the API can restrict its
own responses but cannot restrict a content delivery network.

The infrastructure is fully described as code in eleven Terraform modules across two
environments, and seven GitHub Actions workflows automate linting, unit testing, static security
analysis, container scanning, infrastructure validation, and deployment. The security gate
blocks merges on fixable critical findings rather than merely reporting them.

The automated test suite grew to 117 tests across ten suites. Its practical value was demonstrated
when it caught a defect in the cookie signing implementation that would have prevented every
private video from playing once deployed. A second defect in that same subsystem, which only
appeared once the distribution was live, is described in Section~\ref{sec:challenges}.

\section{Strengths of the System}
\label{sec:strengths}

The strongest property of the design is the separation between the application tier and the media
path. The backend API issues pre-signed URLs, records metadata and answers queries, and never
handles video bytes, which flow directly between the browser, Amazon S3 and CloudFront. The API's
resource consumption is therefore independent of video size and upload concurrency, and a failure
in the processing tier leaves already-published videos playable.

The processing tier costs nothing while idle. Because transcoding runs on AWS Batch over Fargate
rather than on a reserved instance, compute capacity exists only while a job runs.
Section~\ref{sec:cost-analysis} estimates the monthly operating cost of this arrangement at
\$1.51 against \$133.35 for an always-on instance sized for the same workload, a difference that
originates entirely in the elimination of idle capacity.

Content protection is enforced at the edge rather than only in the API. Origin Access Control
keeps the storage bucket private and CloudFront Signed Cookies reject unauthorised requests before
they reach it, which is what makes the private visibility setting meaningful, since an API can
restrict its own responses but cannot restrict a content delivery network. The mechanism was
verified against the deployed system rather than in code alone, as
Section~\ref{sec:challenges} records.

The infrastructure is reproducible. Eleven Terraform modules describe two environments in full,
and seven GitHub Actions workflows automate linting, unit testing, static security analysis,
container and dependency scanning, infrastructure validation and deployment. The security gate
blocks a merge on fixable critical findings rather than merely reporting them, and the scanner
itself is pinned to a release commit rather than to a moving branch.

Finally, the performance claims in this report are measured rather than asserted.
Time-to-First-Frame was measured from six AWS regions and from a consumer connection in Vietnam,
and led to two corrections of the CloudFront price class that reduced the median latency seen from
Vietnam by a factor of 6.6 and the median seen from Brazil by a factor of 32. Transcoding
throughput was measured under one hundred concurrent uploads of real video and cross-validated
with two independently implemented load generators. And the automated suite of 117 tests across
ten suites caught a signing defect that would otherwise have prevented every private video from
playing.

\section{Limitations}
\label{sec:limitations}

Several limitations remain in the delivered system.

The bitrate ladder is fixed at three rungs regardless of content. As the literature reviewed in
Section~\ref{sec:core-models} shows, the optimal ladder is content-dependent, so a fixed ladder
wastes bandwidth on simple content while under-serving complex content. Netflix reported bitrate
reductions between 28\% and 38\% at equivalent perceptual quality by optimising the ladder per
shot \cite{netflix2018dynamic}, which indicates the magnitude of the opportunity forgone.

Each video is transcoded by a single container, without the chunk-based parallelism that Netflix
employs. Processing time therefore scales linearly with video duration, and a long video cannot be
accelerated by adding capacity. Section~\ref{sec:drain-rate} measured roughly 66 seconds of
wall-clock time per job for a three-second clip, which shows how far fixed overhead dominates for
short content; for long content the encoder itself dominates instead, and it is that portion which
the current design cannot parallelise.

Signed Cookies remain valid for two hours after issuance, so switching a video from public to
private leaves viewers holding unexpired cookies with access until those cookies lapse. Immediate
revocation would require rotating the signing key or introducing an edge function that performs a
per-request check. The system also implements no digital rights management: an authorised viewer
can capture the decrypted stream or share their cookies, and preventing this is outside the scope
of the techniques used here.

Rate limiting and view deduplication hold their state in process memory. This is correct for the
current single-instance deployment but would need to move to a shared store such as Redis before
the backend could be scaled horizontally, since each instance would otherwise enforce its own
independent limits. The backend itself remains one process on one instance in a single region and
is the only single point of failure in the system, the transcoding and delivery tiers being
managed services that scale and recover on their own.

Two experimental results are narrower than the headline figures suggest. The comparison between
1 and 4~vCPU allocations rests on a 1~vCPU baseline whose three runs span 56\%, so the estimated
speedup varies between 3.3 and 4.9 times depending on which statistic is compared, and memory was
raised at the same time as the vCPU count, so the experiment does not isolate either variable. The
delivery latency figures were collected from probes running inside AWS regions rather than on
residential connections, so they describe edge proximity under favourable network conditions
rather than what a home viewer would experience.

The two-language interface has one gap of the same kind. A viewer's language choice is held in
the browser and sent to the API on each request, which is sufficient for everything the interface
and the API produce, but the notification e-mails announcing that a video is ready or has failed
are sent from an AWS Batch job long after the request that started it. That job has no request to
read a language from, and the account record stores no language preference, so those two e-mails
are sent in both languages. Storing the preference on the user document would let them be sent in
one, and is the correct fix rather than the one adopted here.

Finally, the CloudFront distributions are provisioned in a second AWS account, because the primary
account was not permitted to create them. The arrangement is invisible to viewers and changes
nothing about the security property being provided, but it splits the infrastructure across two
sets of credentials and carries operational overhead that a verified account would not.

\section{Future Work}
\label{sec:future-work}

Two measurements should still be repeated before the experimental results are considered settled.
The Time-to-First-Frame measurement has been repeated against the live distribution
(Section~\ref{sec:cdn-measurement}) and then from several client locations
(Section~\ref{sec:multiregion-measurement}), which closed that item and, in doing so, exposed a
further limitation that has since been corrected: viewers in South America and Oceania fell
outside the configured price class and were served across an ocean until both distributions were
moved to \texttt{PriceClass\_All}. What remains open is that every probe ran inside an AWS region,
so the figures describe edge proximity under favourable network conditions rather than residential
connections; comparing the two Vietnamese samples suggests the last mile costs roughly a factor of
two, but one pair of observations is not a basis for that claim. The concurrency test re-run with genuine video payloads
(Section~\ref{sec:real-throughput-validation})
confirmed that every job succeeds under load, but it was not instrumented with the same
twenty-five-second polling cadence as the original synthetic run, so a precise drain-rate figure
for real transcoding throughput, as opposed to confirmation that throughput is non-zero and
failure-free, remains to be measured; repeating the same short-clip payload with that polling
protocol in place would close this. The vCPU scaling comparison should also be repeated with more
replications of the 1~vCPU baseline and with memory held constant, since
Section~\ref{sec:pipeline-evaluation} showed the present estimate to be dominated by baseline
variance and confounded by a simultaneous change in memory allocation.

The most valuable functional extension would be per-title bitrate ladder optimisation driven by a
perceptual quality metric such as VMAF. Encoding a short probe of the source at several
operating points and selecting those on the convex hull of the quality-versus-bitrate curve
would reduce delivery bandwidth at unchanged perceived quality, which lowers the largest
remaining cost component of the system.

Chunk-based parallel transcoding would address the linear scaling of processing time. Splitting
the source at keyframe boundaries, encoding the segments concurrently across multiple
containers, and concatenating the results would convert additional capacity into reduced
latency. The main challenge is maintaining rate control consistency across chunk boundaries so
that quality does not fluctuate visibly.

Hardware acceleration offers a further avenue. Research on GPU-enabled serverless platforms for
HEVC encoding \cite{gpuserverless2024} suggests substantial throughput gains, and adopting a
more efficient codec such as HEVC or AV1 alongside H.264 would reduce bandwidth for clients that
support it.

Finally, moving the rate limiting and deduplication state into a shared store, and adding
distributed tracing across the S3, SQS, Lambda, and Batch boundaries, would prepare the system
for horizontal scaling and make failures in the asynchronous pipeline substantially easier to
diagnose.
